\section{Related work}
\label{sec:related_work}

\begin{table}[H]
\centering
\small
\caption{Selected approaches to measurement-error correction. Exactness is understood within each method's stated model and function class.}
\label{tab:related_work}
\renewcommand{\arraystretch}{1.08}
\begin{tabular}{>{\raggedright\arraybackslash}p{0.23\linewidth}>{\raggedright\arraybackslash}p{0.34\linewidth}>{\raggedright\arraybackslash}p{0.35\linewidth}}
\toprule
Work & Main result & Relation to CCSGD \\
\midrule
Nakamura (1990) \cite{nakamura1990} & Corrected estimating equations & Foundational score-correction principle. \\
Novick--Stefanski (2002) \cite{novick2002} & Exact Gaussian complex Monte Carlo corrected scores and batch $M$-estimation & Our Gaussian correction uses their identity; we study its per-sample use in SGD. \\
Stefanski (1989) \cite{stefanski1989} & Analyticity restriction for globally unbiased Gaussian correction & Delimits exact correction for scores with complex singularities. \\
Konstantopoulos--Patie--Sarkar (2024) \cite{konstantopoulos2024} & van Dantzig pairs and distributional constructions & Supplies candidate conjugate perturbations for non-Gaussian noise; score integrability remains to be checked. \\
Woolsey--Huang (2026) \cite{woolseyhuang2026} & Statistical estimation with non-Gaussian error through finite-order bias cancellation & CCSGD gives an explicit stochastic-gradient update and optimization guarantees; exact correction for admissible non-Gaussian noise is discussed in Further work. \\
Cook--Stefanski (1994); Spicker et al. (2024) \cite{cookstefanski1994,spicker2024} & Parametric and nonparametric SIMEX & Correct by extrapolation rather than an exact per-sample gradient identity. \\
Belomestny--Goldenshluger (2021) \cite{belomestnygoldenshluger2021} & Density deconvolution allowing zeros of the error characteristic function & Related inverse problem; its target is density estimation. \\
Loh--Wainwright (2012) \cite{lohwainwright2012} & Corrected high-dimensional sparse linear regression with a convergent projected-gradient algorithm & Addresses linear regression; CCSGD treats scores of more general analytic losses in an online stochastic update. \\
Carroll et al. (2006) \cite{carroll2006} & Regression calibration and broader measurement-error methodology & Context for plug-in correction. \\
\bottomrule
\end{tabular}
\end{table}

Nakamura's corrected-score principle and the complex Monte Carlo
construction of Novick and Stefanski \cite{nakamura1990,novick2002}
provide the statistical foundation. CCSGD studies the associated online
gradient recursion and derives its noise conditions before applying
classical optimization theory. Loh and Wainwright
\cite{lohwainwright2012} already provide a computationally tractable
corrected optimization method for sparse linear regression; our analysis
allows more general analytic scores under explicit growth conditions.
Woolsey and Huang \cite{woolseyhuang2026} develop non-Gaussian
corrections for statistical estimation by cancelling finitely many bias
terms, which generally leaves a residual bias. Our focus is stochastic
optimization: CCSGD gives a directly implementable learning rule using
one noisy observation and one simulated complex perturbation per update,
with convergence rates and asymptotic guarantees. Further work discusses
exact correction for non-Gaussian noise admitting a suitable conjugate
perturbation, including uniform and symmetric two-point noise, under the
required score-integrability conditions. In our experiments, the
Woolsey--Huang baseline is their deterministic correction in Eq.~(5).

\subsection*{Our contributions}
\begin{enumerate}
\item \emph{An explicit optimization method.}
CCSGD turns complex correction into a directly implementable stochastic
optimization algorithm for a broad class of analytic losses with noisy
inputs. Each update evaluates the score at one noisy observation plus
one simulated complex perturbation. Experiments demonstrate convergence
to the true parameter in nonlinear-feature regression and improved
prediction over naive training in neural-network regression.
\item \emph{Exact correction beyond Gaussian noise.}
The conjugate-variable criterion in our companion work
\cite{conjugatecriterion2026} extends the correction principle to
admissible non-Gaussian noise, including uniform and symmetric two-point
distributions. Further work describes these constructions and the
integrability conditions required for exact unbiasedness; the main
optimization theorems here concern Gaussian measurement error.
\item \emph{Per-step unbiasedness and asymptotic guarantees.}
The corrected gradient is exactly unbiased at every iteration, rather
than merely cancelling finitely many bias terms. We derive the classical
SGD relative variance bound from input growth and sampling conditions,
then establish convex and nonconvex guarantees. Under strong convexity,
CCSGD converges to the clean optimizer $\theta^*$ almost surely and in
mean square, with explicit MSE rates; under the stated local regularity
conditions, both the last iterate and the Polyak--Ruppert average are
asymptotically normal.
\end{enumerate}
